\documentclass[a4paper]{article}

\usepackage{graphicx}
\usepackage{amsmath,amssymb}
\usepackage{minted}
\usepackage{multirow}
\usepackage{float}
\usepackage{todonotes}
\usepackage{forest}
\usepackage{xstring}
\usepackage{xcolor}
\usepackage[most]{tcolorbox}
\usepackage{xparse}
\usepackage{natbib}

\usetikzlibrary{graphs,graphdrawing}
\usegdlibrary{layered}

% for external graphviz
\input{/home/donkarlo/Dropbox/repo/nd_ai_project/src/nd_ai/knoweldge_representation/language/formal/computer/programming/tex/latex/code_clips/external_graphviz.tex}

% for tags
\input{/home/donkarlo/Dropbox/repo/nd_ai_project/src/nd_ai/knoweldge_representation/language/formal/computer/programming/tex/latex/parts/control_sequence/kind/semtag/semtag.tex}

% for links
\input{/home/donkarlo/Dropbox/repo/nd_ai_project/src/nd_ai/knoweldge_representation/language/formal/computer/programming/tex/latex/code_clips/seehere.tex}

\title{Feature structure (trace)}
\date{}
\author{Mohammad Rahmani}

\begin{document}
    \maketitle
    \seeherefooter{https://en.wikipedia.org/wiki/Feature_structure}
    \cite{kay-1984-functional-unification-grammar} defines ``Feature structure`` as a functional description of a linguistic object. It is built from features, where each feature consists of an attribute and an associated value. A value may be atomic, or it may itself be another functional description, so the structure can be nested. Functional descriptions are combined by unification when they are compatible.
    
    Feature structures \cite{kasper-1986-a-logical-semantics-for-feature-structures} represent linguistic objects as sets of attribute--value pairs. Kasper and Rounds formalize feature-structure descriptions as logical formulas interpreted by directed graphs.
    Their model explains unification and shows how disjunction and non-local path values increase computational complexity.
    
    \bibliographystyle{plainnat}
    \bibliography{/home/donkarlo/Dropbox/repo/nd_shared_research/bibliography/bibliography.bib}
\end{document}